\section{Counting the coefficients of a formal circuit}\label{sec:parse}

A circuit that computes $P_n^*$ exactly supplies a coefficient for each
input to a signed processor. A border circuit supplies a point in the
closure of a family of coefficient vectors. To use such a point in the
clock argument, we first replace every processor entry by a polynomial
in a common set of scalar parameters $X$. Each path of the resulting
contraction contributes a monomial in $X$. Counting the paths that
contribute a given monomial will let us apply the compression of
Theorem~\ref{thm:gap}.

Assume for contradiction that $\bsz_F^{[n]}(P_n^*)\le m_0$, where
\[
 m_0=\lfloor 2^{n^\sigma}\rfloor,\qquad
 R=2^{\lceil n^\rho\rceil},\qquad M=m_0+n+w .
\]
Here $0<\sigma<\rho$ and $w$ is the word length from
Section~\ref{sec:gap}. Thus $M^C=o(R)$ for every fixed $C$.
The coefficient vector of $P_n^*$ is defined over the prime field.
The border hypothesis holds over its algebraic closure, since the
kernels of the circuit coefficient maps commute with field extension,
as in Section~\ref{sec:comparison}.

\subsection{One table for all circuits of a given size}
A coefficient query at $x\in\{0,1\}^n$ asks for the coefficient of
$\prod_i Z_{i,x_i}$. We parametrize the entire degree-$n$ output,
including the coefficients outside this set. The recursive form of the
parametrization will allow all $R$ processors to answer their queries
within the larger clock.

\begin{lemma}[Parse tables]\label{lem:parse-tables}
For every $n\ge1$ and $s\ge 1$ there is an integral polynomial map
$Y(X)$ from $h=(n+s)^{O(1)}$ scalar parameters to the full space of
degree-$n$ coefficients in $2n$ variables. In every field, the image
contains the degree-$n$ parts of all circuits of size at most $s$.
Each coordinate of $Y$ has degree at most $2n-1$ in $X$, and its
coefficient $\ell^1$-norm has logarithm $O(n\log(n+s+2))$ when nonzero.
\end{lemma}
\begin{proof}
Reserve $G=O(n+s)$ gate slots, including the used inputs and a
distinguished output. For each slot $g$ define $T_{g,1}$ to be a
linear combination of the $2n$ variables, with a separate parameter
for every coefficient. Recursively define
\[
 T_{g,d}=\sum_{e=1}^{d-1}\sum_{a,b=1}^{G}
     X_{g,d,e,a,b}\,T_{a,e}T_{b,d-e}\qquad (2\le d\le n).
\]
Take the coefficients of $T_{\mathrm{out},n}$ for the coordinates of
$Y$. The table has $O(nG+n^2G^3)$ parameters.

To encode a particular circuit, evaluate its degree-zero gate
components first. At each positive degree $d$, addition, subtraction
and multiplication by a degree-zero component give linear combinations
of earlier degree-$d$ components. Expand these linear dependencies in
topological order. At degree one the resulting expressions are linear
forms in the input variables. At degree $d\ge2$ they are linear
combinations of products of two positive-degree components with degrees
summing to $d$. Their coefficients give the required parameter values.

Expand the recursion before specializing any parameter. Each summand
corresponds to an ordered full binary tree with $n$ leaves. At each
internal node choose a split and two child slots. At each leaf choose
an input variable. A tree therefore contributes a product of $2n-1$
parameters. There are fewer than $4^n$ tree shapes, each with at most
$(n+s+2)^{O(n)}$ labelings. Collecting terms gives the asserted norm
bound, including multiplicities when two trees yield the same monomial.
\end{proof}

\subsection{Counting a coefficient of the contraction}
Fix the private machine of Theorem~\ref{thm:gap}(2) and the
$R$-processor wiring of Theorem~\ref{thm:gap}(1). The public description
$D$, including its code, uses at most $R/10$ active words. Set all
other description words to zero. Write $F_D(y)$ for the contraction
with these pins when each processor with input $x$ contributes the
coordinate $y_x$ indexed by $\prod_i Z_{i,x_i}$. In particular
$F_D(c)=Z(D)$, the genuine completed success count, for the full
coefficient vector $c$ of $P_n^*$.

At size $s=m_0$ the substituted contraction is
\[
 A_D(X)=F_D(Y(X)),\qquad \deg_X A_D\le L:=R(2n-1).
\]
Indeed, each assignment of the linked bits yields a product of $R$
coordinates of $Y$. All coefficients of $A_D$ are nonnegative integers.
Its value at the boundary will be recovered only after these terms
are summed.

\begin{lemma}[Coefficient counts]\label{lem:substitution}
Given $D$ and an exponent vector $\mu\in\mathbb Z_{\ge0}^{h}$ with
$|\mu|\le L$, a binary-guess program without private jobs at cap
$T_+$ has exactly $[X^\mu]A_D$ successful paths. The encoding of
$\mu$ uses $M^{O(1)}$ bits. The assertion holds for sufficiently small
$\rho$ with $0<\sigma<\rho$.
\end{lemma}

\begin{proof}
Generate the pins and wiring from $D$ and guess the linked bit
values. At each of the resulting $R$ inputs $x$, guess an ordered
tree from Lemma~\ref{lem:parse-tables}. Encode its $2n-1$ nodes in
preorder. Each record specifies a gate slot and degree. A leaf gives
its input variable, and an internal node gives its degree split and
ordered child slots. Check all ranges, the root slot and degree, and
require unused record fields to be zero.

A degree-$d$ node at position $u$ occupies $[u,u+2d-2]$ within
its block. For a split $d=e+(d-e)$ check the child records at
$u+1$ and $u+2e$. Their intervals partition the rest of the parent's
interval. Induction from the root therefore verifies every record
exactly once. Tally the leaf variables and require the monomial to be
$\prod_i Z_{i,x_i}$. Across all trees, tally the uses of every
parameter in a common random access array and require the totals
to equal $\mu$. A successful path specifies one wiring assignment and
one labeled ordered tree at each processor, exactly the multiplicity
with which that choice occurs in $[X^\mu]A_D$.

Each count, slot label or address fits in $O(w)$ bits. Initializing and
comparing the parameter totals costs $M^{O(1)}$, since the table has
$m_0^{O(1)}+n^{O(1)}$ parameters and $\log(L+1)=O(w)$. Generating the
wiring, checking the records and tallying the leaf variables takes
$O(Rn\operatorname{poly}(w))$ ticks. The exponent on $w$ is fixed, so
decreasing $\rho$ makes this term fit $T_+$, while $M^{O(1)}=o(R)$.
\end{proof}

Theorem~\ref{thm:gap}(2) now gives a smaller-clock simulation of each
coefficient count. Its compressor parameters have fixed values. The
parse-table parameters $X$ still have to account for a boundary point
$c$ of the image of $Y$, possibly with no assignment $X$ for which
$Y(X)=c$.

Sections~\ref{sec:descent} and~\ref{sec:cartier} recover this boundary
value through matrices over a prime field. One matrix $W_\mu$ will
represent each monomial $X^\mu$. The same nonzero vector $\Omega$ will
satisfy
\[
 \biggl(\sum_\mu ([X^\mu]A_D)\,W_\mu\biggr)\Omega=Z(D)\Omega
\]
for every description $D$. The computation can guess $\mu$ and count
one matrix entry using only polynomially many bits in $M$. In
Section~\ref{sec:border} it will supply one index as part of its own
description and guess the other. The vector $\Omega$ enters there
only through the displayed identity and the choice of one index at
which it is nonzero.
